\section{Conclusion}
\label{sec:conclusion}

Standard IEC/IEEE~60255-118-1 tests do not fully characterize frequency
estimator behavior in converter-dominated grids, where simultaneous phase
discontinuities, high RoCoF, and harmonic distortion expose failure modes
that isolated single-disturbance tests miss. A benchmark of twelve estimators
across five algorithmic families and five stress scenarios yields three
principal observations.

\begin{enumerate}

\item \textbf{Latency--Robustness Trade-off:}
Window-based methods achieve strong steady-state accuracy, but their
behavior differs across scenario types. IpDFT is better suited to
metering-oriented operation, whereas TFT reduces observation delay at
the cost of higher noise sensitivity under composite disturbances.
Loop-based estimators (SRF-PLL, SOGI-FLL) are computationally efficient
and handle isolated phase jumps well, but exhibit elevated trip-risk
under the composite Multi-Event sequence
($T_{\mathrm{trip}}\!=\!0.353$\,s for SRF-PLL; Table~\ref{tab:risk_metrics}).

\item \textbf{Model-Based Filters:}
Within the model-based family, the RA-EKF (augmented with an explicit
$\dot{\omega}$ state, Innovation-Driven Covariance Scaling, and Event Gating)
achieves $T_{\mathrm{trip}}\!=\!0$ under Composite Islanding vs.\
$T_{\mathrm{trip}}\!=\!0.120$\,s for the standard EKF ($N\!=\!30$ Monte Carlo),
at $88.5\,\mu$s/sample (Python baseline; C/HIL validation remains future work).
The standard EKF and UKF offer lower cost but higher trip-risk under composite
disturbances; all three model-based filters require careful tuning for
IBR-specific disturbance types.

\item \textbf{Limits of Standard Compliance Tests:}
Several estimators passing IEC/IEEE~60255-118-1 step and ramp tests
exhibited degraded performance under composite multi-disturbance
scenarios, suggesting that standard compliance tests may not fully
capture estimator resilience in IBR-dominated conditions.

\end{enumerate}

Data-driven approaches show variable cost: PI-GRU requires
${\approx}154$\,ms/sample (Python/PyTorch CPU baseline), prohibitive for
embedded protection without GPU acceleration, and did not outperform
scenario-optimized classical methods in any of the five evaluated scenarios.
This benchmark is simulation-based; hardware-in-the-loop validation remains
a next step.
The results suggest that IEC/IEEE~60255-118-1 compliance tests do not
fully characterize estimator resilience in IBR-dominated grids; a
\emph{Composite Phase-Step with Harmonic Distortion} scenario---analogous
to Scenarios~D--E---with probabilistic compliance bounds
($\mathrm{PCB} = \mu_{|e|} + 3\sigma_{|e|}$) may better assess relay
security. Simulation code and hyperparameter grids are available from
the authors for reproducibility.